\documentclass[12pt]{article}
\usepackage[utf8]{inputenc}
\usepackage{listings}
\usepackage{xcolor}
\usepackage{geometry}
\usepackage{parskip}

\geometry{margin=2.5cm}

\definecolor{codebg}{RGB}{245,245,245}
\definecolor{keyword}{RGB}{0,112,0}
\definecolor{comment}{RGB}{128,128,128}
\definecolor{string}{RGB}{180,0,0}

\lstdefinestyle{pythonstyle}{
    backgroundcolor=\color{codebg},
    basicstyle=\ttfamily\small,
    keywordstyle=\color{keyword}\bfseries,
    stringstyle=\color{string},
    commentstyle=\color{comment}\itshape,
    language=Python,
    frame=single,
    rulecolor=\color{gray!40},
    breaklines=true,
    showstringspaces=false,
    tabsize=4,
    captionpos=b,
}

\title{\textbf{Wind Turbine Blade Fault Classification\\Using a Simple CNN}\\[0.5em]
\large Step 11 --- Plot Accuracy and Loss Graphs}
\author{Amreen Batool}
\date{}

\begin{document}
\maketitle

\section*{Step 11 --- Plot Accuracy and Loss Graphs}

\subsection*{Accuracy Plot}

\begin{lstlisting}[style=pythonstyle]
plt.figure(figsize=(8, 5))
plt.plot(history.history['accuracy'],     marker='o', label='Training Accuracy')
plt.plot(history.history['val_accuracy'], marker='s', label='Validation Accuracy')
plt.xlabel('Epoch')
plt.ylabel('Accuracy')
plt.title('Training and Validation Accuracy')
plt.legend()
plt.grid(True)
plt.show()
\end{lstlisting}

\subsection*{Loss Plot}

\begin{lstlisting}[style=pythonstyle]
plt.figure(figsize=(8, 5))
plt.plot(history.history['loss'],     marker='o', label='Training Loss')
plt.plot(history.history['val_loss'], marker='s', label='Validation Loss')
plt.xlabel('Epoch')
plt.ylabel('Loss')
plt.title('Training and Validation Loss')
plt.legend()
plt.grid(True)
plt.show()
\end{lstlisting}

\subsection*{Explanation}

These two graphs together show how the model learned over time:

\begin{itemize}
    \item \textbf{Accuracy graph} --- how training and validation accuracy improved across epochs. Both curves rising together is a good sign.
    \item \textbf{Loss graph} --- how the error decreased across epochs. Both curves falling together is a good sign.
    \item If training accuracy is much higher than validation accuracy, this suggests \textit{overfitting}.
    \item If both are low and similar, the model may be \textit{underfitting}.
\end{itemize}

\subsection*{Figures to Include}

\begin{itemize}
    \item \textbf{Figure 2: Training and Validation Accuracy}
    \item \textbf{Figure 3: Training and Validation Loss}
\end{itemize}

\end{document}
